% GENERATED FILE. Edit canonical lessons, then run npm run generate:books.
% canonical-source-hash: fnv1a64:fb471b99e7aa7959
% canonical-lessons: TE-C228-pita, TE-C228-telu, TE-C228-gabbilam, TE-C228-gudu, TE-C228-toka

\chapter{Crab, Scorpion, Bat, Nest, Tail}
\label{ch:te-crab-scorpion-bat-nest-tail}
\hlchaptermodality{228}

\begin{chapteropening}
\textbf{By the end of this chapter:} \emph{I can say a crab, a scorpion, a bat, a nest and a tail.}

Complete the last lesson of chapter 228: I can say a crab, a scorpion, a bat, a nest and a tail.
\end{chapteropening}

\section[\texorpdfstring{\te{పీత} (pīta)}{పీత (pīta)}]{\te{పీత} (pīta) — a crab}

\label{lesson:TE-C228-pita}



\begin{quote}
Before the new one: say the Telugu for an enemy, then the Telugu for a bridegroom.
\end{quote}

\subsection*{You'll want to know: \te{పీత}}
\textbf{\te{పీత}} — \emph{pīta} — \textquotedblleft{}a crab\textquotedblright{}.

\begin{cousinweb}
The origin of \textbf{\te{పీత}} (\emph{pīta}) is uncertain. DEDR and CDIAL list no word of this shape meaning \textquotedblleft{}crab\textquotedblright{}; CDIAL's \emph{pīta}, \textquotedblleft{}drunk\textquotedblright{} and \textquotedblleft{}yellow\textquotedblright{}, are different words.
\end{cousinweb}

\begin{grammarlens}[title={What you've built}]
One more word for this chapter.
\end{grammarlens}

\subsection*{Your turn}
\begin{itemize}
\raggedright
  \item \emph{Say it:} \emph{pīta}
  \item \emph{Say it:} \emph{pīta}, once more
  \item \emph{Say it:} say \emph{pīta}
  \item \emph{Recall:} say the Telugu for an enemy, then the Telugu for a bridegroom, then say \emph{pīta} again
\end{itemize}

\begin{tcolorbox}[breakable,colback=teal!4,colframe=teal!35!black,title={Before you move on}]
What does \te{పీత} mean? (\textquotedblleft{}A crab\textquotedblright{}.) Say it once more.
\end{tcolorbox}
\section[\texorpdfstring{\te{తేలు} (tēlu)}{తేలు (tēlu)}]{\te{తేలు} (tēlu) — a scorpion}

\label{lesson:TE-C228-telu}



\begin{quote}
Before the new one: say the Telugu for a bridegroom, then the Telugu for a crab.
\end{quote}

\subsection*{You'll want to know: \te{తేలు}}
\textbf{\te{తేలు}} — \emph{tēlu} — \textquotedblleft{}a scorpion\textquotedblright{}.

\begin{cousinweb}
\textbf{\te{తేలు}} (\emph{tēlu}) is native Dravidian, one word with Tamil \emph{tēḷ} and Malayalam \emph{tēḷ} (DEDR 3470). DEDR lists Telugu \emph{tēlu}, \textquotedblleft{}scorpion\textquotedblright{}.
\end{cousinweb}

\begin{grammarlens}[title={What you've built}]
One more word for this chapter.
\end{grammarlens}

\subsection*{Your turn}
\begin{itemize}
\raggedright
  \item \emph{Say it:} \emph{tēlu}
  \item \emph{Say it:} \emph{tēlu}, once more
  \item \emph{Say it:} say \emph{tēlu}
  \item \emph{Recall:} say the Telugu for a bridegroom, then the Telugu for a crab, then say \emph{tēlu} again
\end{itemize}

\begin{tcolorbox}[breakable,colback=teal!4,colframe=teal!35!black,title={Before you move on}]
What does \te{తేలు} mean? (\textquotedblleft{}A scorpion\textquotedblright{}.) Say it once more.
\end{tcolorbox}
\section[\texorpdfstring{\te{గబ్బిలం} (gabbilaṁ)}{గబ్బిలం (gabbilaṁ)}]{\te{గబ్బిలం} (gabbilaṁ) — a bat (animal)}

\label{lesson:TE-C228-gabbilam}



\begin{quote}
Before the new one: say the Telugu for a crab, then the Telugu for a scorpion.
\end{quote}

\subsection*{You'll want to know: \te{గబ్బిలం}}
\textbf{\te{గబ్బిలం}} — \emph{gabbilaṁ} — \textquotedblleft{}a bat (animal)\textquotedblright{}.

\begin{cousinweb}
\textbf{\te{గబ్బిలం}} (\emph{gabbilaṁ}) is native Dravidian. DEDR 1216 lists \emph{gabbilamu}, \textquotedblleft{}bat\textquotedblright{}, with Kannada \emph{gabbilāyi} and \emph{kapaṭe}, \textquotedblleft{}bat\textquotedblright{}.
\end{cousinweb}

\begin{grammarlens}[title={What you've built}]
One more word for this chapter.
\end{grammarlens}

\subsection*{Your turn}
\begin{itemize}
\raggedright
  \item \emph{Say it:} \emph{gabbilaṁ}
  \item \emph{Say it:} \emph{gabbilaṁ}, once more
  \item \emph{Say it:} say \emph{gabbilaṁ}
  \item \emph{Recall:} say the Telugu for a crab, then the Telugu for a scorpion, then say \emph{gabbilaṁ} again
\end{itemize}

\begin{tcolorbox}[breakable,colback=teal!4,colframe=teal!35!black,title={Before you move on}]
What does \te{గబ్బిలం} mean? (\textquotedblleft{}A bat (animal)\textquotedblright{}.) Say it once more.
\end{tcolorbox}
\section[\texorpdfstring{\te{గూడు} (gūḍu)}{గూడు (gūḍu)}]{\te{గూడు} (gūḍu) — a nest}

\label{lesson:TE-C228-gudu}



\begin{quote}
Before the new one: say the Telugu for a scorpion, then the Telugu for a bat.
\end{quote}

\subsection*{You'll want to know: \te{గూడు}}
\textbf{\te{గూడు}} — \emph{gūḍu} — \textquotedblleft{}a nest\textquotedblright{}.

\begin{cousinweb}
\textbf{\te{గూడు}} (\emph{gūḍu}) is native Dravidian. DEDR 1883 lists a nasalised \emph{gūḍu}, \textquotedblleft{}nest, cage\textquotedblright{}, with Tamil \emph{kūṭu} and Kannada \emph{gūḍu}, \textquotedblleft{}nest\textquotedblright{}. DEDR compares Sanskrit \emph{kūṭa}, \textquotedblleft{}trap for catching deer\textquotedblright{} (CDIAL 3397).
\end{cousinweb}

\begin{grammarlens}[title={What you've built}]
One more word for this chapter.
\end{grammarlens}

\subsection*{Your turn}
\begin{itemize}
\raggedright
  \item \emph{Say it:} \emph{gūḍu}
  \item \emph{Say it:} \emph{gūḍu}, once more
  \item \emph{Say it:} say \emph{gūḍu}
  \item \emph{Recall:} say the Telugu for a scorpion, then the Telugu for a bat, then say \emph{gūḍu} again
\end{itemize}

\begin{tcolorbox}[breakable,colback=teal!4,colframe=teal!35!black,title={Before you move on}]
What does \te{గూడు} mean? (\textquotedblleft{}A nest\textquotedblright{}.) Say it once more.
\end{tcolorbox}
\section[\texorpdfstring{\te{తోక} (tōka)}{తోక (tōka)}]{\te{తోక} (tōka) — a tail}

\label{lesson:TE-C228-toka}



\begin{quote}
Before the new one: say the Telugu for a bat, then the Telugu for a nest.
\end{quote}

\subsection*{You'll want to know: \te{తోక}}
\textbf{\te{తోక}} — \emph{tōka} — \textquotedblleft{}a tail\textquotedblright{}.

\begin{cousinweb}
\textbf{\te{తోక}} (\emph{tōka}) is native Dravidian. DEDR 3538 lists a nasalised \emph{tōka}, \textquotedblleft{}tail\textquotedblright{}, with Tamil \emph{tōkai} and Kannada \emph{tōke}, \textquotedblleft{}tail\textquotedblright{}.
\end{cousinweb}

\begin{grammarlens}[title={What you've built}]
One more word for this chapter.
\end{grammarlens}

\subsection*{Your turn}
\begin{itemize}
\raggedright
  \item \emph{Say it:} \emph{tōka}
  \item \emph{Say it:} \emph{tōka}, once more
  \item \emph{Say it:} say \emph{tōka}
  \item \emph{Recall:} say the Telugu for a bat, then the Telugu for a nest, then say \emph{tōka} again
\end{itemize}

\begin{tcolorbox}[breakable,colback=teal!4,colframe=teal!35!black,title={Before you move on}]
What does \te{తోక} mean? (\textquotedblleft{}A tail\textquotedblright{}.) Say it once more.
\end{tcolorbox}
